\scanpage{8}\mspage{7}
\[
\begin{aligned}
\leq,\quad\ll,\quad\preccurlyeq&\qquad\text{relations d'ordre}\\
&\qquad\text{\unsure{dans|de}}\ \unsure{\mathcal{F}|\mathfrak{F}}
\end{aligned}\tag*{(1.10).}
\]
et les relations induites dans $\mathcal{M}$
\[
\leq,\quad\ll\quad\text{dans }\mathcal{M}\tag*{(1.11)}
\]
% slanted
\margin{\uline{NB} les relations d'ordre induites par $\leq$, $\ll$ sur $\mathcal{L}$ sont discrètes}
[nb la relation d'ordre $\preccurlyeq$ induite dans $\mathcal{M}$ est moins importante, comme structure spécifique dans $\mathcal{M}$].

% slanted
\margin{\uline{NB} $G\leq F\Longrightarrow G\ll F$. une sous-figure de $F$ raffine $F$}
\begin{center}
\begin{tabular}{@{}l@{\qquad}p{.64\linewidth}@{}}
$\deleted{F}G\deleted{\unsure{<|\ll}}\leq\deleted{G}F$ & $\deleted{F}G$ est dit une \uline{sous}-\uline{figure} de $\deleted{G}F$\\
& (Si $\deleted{F}G\inserted{=X}\in\mathcal{M}$, c'est donc une \uline{strate} de $\deleted{G}F$)\\[.45em]
$G\ll F$ & $G$ est dit un \uline{raffinement} de $F$\\
& (Si $G=\unsure{x|X}\in\mathcal{M}$, c'est donc un élément de la grande o\unsure{m|\deleted{o}m}bre de $F$)\\[.45em]
$G\preccurlyeq F$ & $G$ est dit une \uline{subdivision} de $F$,\\
& on dit que la figure $G$ subdivise la figure $F$.
\end{tabular}
\end{center}

Je note enfin deu\deleted{\unsure{x}}x autres relations importantes \deleted{dans} entre figures:
\[
\left\{\begin{aligned}
F\text{ et }G&\quad\text{\uline{compatibles}}\\
F\text{ et }G&\quad\text{\uline{disjoints}}\qquad,
\end{aligned}\right.\tag*{(1.12)}
\]
qui l'une et l'autre peuvent s'exprimer en termes de la relation d'ordre $\leq$:

$F$ et $G$ compatibles: ce sont des sous figures d'une $\overline m$ figure, i..e $\{F,G\}$ majoré pour $\leq$;

$F$ et $G$ disjoint: $F$ et $G$ compatibles, et leur
